\chapter{语音任务}
\label{chap:nlp-speech-tasks}

用户说“读书会是四月十八日吗”，系统即使写对了每个字，也可能在用户尚未说完时抢答，或者把正确的日期读成“四月十七日”。第\ref{chap:nlp-question-answering-dialogue}章已经区分回答依据与对话状态；声音接口又带来两个独立问题：连续的声学观测怎样对应离散语言单位，正在发生的交流应在何时接收、确认和输出。识别、口语理解、合成、翻译与实时对话分别交付文字、语义、声音、跨语言结果和完整交互，不能用其中一项的正确代替其余各项。

本章的声音情境、概率、时长和运行记录均为教学构造，不表示实际录音、模型训练或性能测量。共同材料仍是2025年4月15日发布的机器学习读书会通知：活动\code{ML-0418}原定4月17日，改至4月18日（周五）14:00，地点为海棠楼201；校内师生可直接参加，校外来宾须在4月16日18:00前提交报名信息。口语可能把“14:00”读成“下午两点”，也可能在日期之间插入停顿。转写规范需要规定是否保留口头形式，语义核验则需要保持同一时刻和严格截止条件。

先把声音中的时间与程序运行的时间分开。设录音中某个单位的结束位置为$a_i$，在线系统开始接收这段录音的墙钟时刻为$t_0$；实时匀速送入时，该单位最早完整可见的时刻为$t_0+a_i$。文本单位首次不可撤回地提交于$c_i$，则其提交延迟为$c_i-(t_0+a_i)$。整轮用户结束于$t_{\mathrm{end}}$，最终转写、首段回复开始播放和业务完成分别发生于$t_{\mathrm{asr}}$、$t_{\mathrm{play}}$、$t_{\mathrm{done}}$。这三个时刻可以相差很远；离线文件快速解码也不能冒充实时匀速输入。

\begin{table*}[htbp]
  \centering
  \begin{tabularx}{\linewidth}{@{}p{15mm}p{28mm}X X p{36mm}@{}}
    \toprule
    任务 & 输入与输出 & 正确性及局部错误 & 端到端要求 & 时间起止点\\
    \midrule
    识别 & 声音到文字、时间区间 & 转写编辑距离；另查边界与说话者 & 该说的没有遗漏，未说的不新增 & 单位声学结束到稳定提交；整轮结束到最终转写\\
    口语理解 & 声音到意图、槽位或事实 & 日期、否定、更正及完整记录 & 当前有效要求正确，来源可追踪 & 所需证据可见到语义记录确认\\
    合成 & 文字及声音条件到波形 & 数量、读音、遗漏；自然度另评 & 听者获得规定内容和表达 & 请求到首段可播放；请求到全部声音生成\\
    翻译 & 源声音到目标文字或声音 & 原意、实体与额外新增内容 & 译文正确且实际声音也正确 & 源语证据可见到目标提交；另计首段播放和全段播放结束\\
    实时对话 & 双方声音、历史、工具到交互 & 抢话、漏听、状态错误、重复动作 & 合理轮次、可听回复及业务完成证据 & 用户结束到首段播放；打断到停止；请求到业务完成\\
    \bottomrule
  \end{tabularx}
  \caption{语音任务的产出、误差与时间口径。同一日期询问可以转写正确却回复过早。表中的起止点须由同一时钟或明确换算得到；波形生成结束、播放结束和业务完成不是同一个事件。}
  \label{tab:nlp-speech-contracts}
\end{table*}
\FloatBarrier

\section{语音识别与时间对齐}
\label{sec:nlp-speech-asr}

\term{自动语音识别}（automatic speech recognition, ASR）把声音转换为规定形式的文字。最小输出是符号序列；分段、标点、说话者和时间戳都是额外产物。例如“四月十八日下午两点”可以被正确转写，却错误地标在另一位说话者的区间，或者把起点提前到静音段。表\ref{tab:nlp-speech-contracts}因此把文字与时间、身份分别验收。

\subsection{声音输入与转写监督}
\label{sec:nlp-speech-input}

一个识别样本包含波形、采样率、语言或允许语种、转写规范及来源标识；有标注时再附说话者和时间区间。字符、子词和词作为输出单位时会形成不同词表，空格、语气词、数字与标点的取舍也影响错误率。训练前应按原始录音、会话或说话者分组划分数据，再截段和增强，不能把同一录音的相邻切片同时放进训练集和测试集。

人工转写提供经过核查的配对监督；网络字幕等弱监督可能包含错位、删节和非语音文字；无转写录音则可先学习声音表征，再用配对样本训练识别头。自监督与弱监督的通用机制见第\ref{chap:machine-learning-paradigms}章。对同一句通知，口音改变发音实现，背景噪声改变观测，停顿改变时序；这些变化不必改变规范文字，但其中一部分若已不可听清，沿用完整标签就需要额外审查。

\subsubsection{声学特征：log-Mel与MFCC}
\label{sec:nlp-speech-features}

波形采样记录瞬时振幅，直接观察它不容易区分某个音的频率成分。短时分析把信号切成相互重叠的小窗，在每窗内近似观察局部频率结构。设单声道波形为$x_n$，采样率$f_s$，窗长$w$个采样，帧移$h$个采样。需要统一采样率时，应先作抗混叠滤波再重采样；仅改文件头会同时改变读出的速度和音高。

对第$i$帧乘窗函数$g_n$，补零到变换长度$k_{\mathrm{fft}}\geq w$，再计算\term{短时傅里叶变换}（short-time Fourier transform, STFT）：
\begin{equation}
 S_{i,k}=\sum_{n=0}^{w-1}x_{ih+n}g_n
       \exp\!\left(-\mathrm{i}\frac{2\pi kn}{k_{\mathrm{fft}}}\right),
 \qquad P_{i,k}=\frac{|S_{i,k}|^2}{k_{\mathrm{fft}}}.
 \label{eq:nlp-speech-stft}
\end{equation}
这里$i$从0开始，$\mathrm{i}$为虚数单位；实信号只保留$k=0,\ldots,k_{\mathrm{fft}}/2$。功率谱$P_{i,k}$经过$m$个梅尔滤波器$H_{j,k}$汇总，再取对数，得到\term{对数梅尔滤波器组特征}（log-Mel）：用较少的频带描述每帧的能量分布，
\begin{equation}
 \bm{X}_{i,j}=\log\!\left(\epsilon+
                   \sum_k H_{j,k}P_{i,k}\right),
 \quad j=1,\ldots,m.
 \label{eq:nlp-speech-logmel}
\end{equation}
$\epsilon>0$避免对零取对数，滤波器按预先选定的梅尔频率尺度布置。滤波器定义、对数底数和归一化方式都应随模型保存，不同实现的“80维特征”不一定数值相同。

给一秒、16\,kHz的教学波形，共16000个采样。取25\,ms窗长$w=400$、10\,ms帧移$h=160$、$k_{\mathrm{fft}}=512$、$m=80$，不做居中补边，也丢弃不足一窗的尾部，则帧数为
\begin{equation}
 n_{\mathrm{frame}}=
 1+\left\lfloor\frac{16000-400}{160}\right\rfloor=98.
 \label{eq:nlp-speech-frame-count}
\end{equation}
因此STFT单边谱为$98\times257$，log-Mel矩阵为$98\times80$。最后一窗从0.970\,s到0.995\,s，最后5\,ms未进入完整窗。若改用20\,ms帧移，得到49帧；保持10\,ms帧移、改用40\,ms窗，则得到97帧。更长的窗提供更长的频率观察区间，但会混合更长时间的变化；补零只是细化离散频率采样，不能凭空增加观测分辨率。居中补边实现可能得到不同帧数，必须注明边界约定。

\term{梅尔倒谱系数}（Mel-frequency cepstral coefficients, MFCC）在滤波器组的对数能量上再做离散余弦变换。保留前$q$项时，可写
$c_{i,r}=\sum_{j=0}^{m-1}\bm{X}_{i,j+1}\cos[\pi r(j+1/2)/m]$，
$r=0,\ldots,q-1$，归一化系数依实现约定。取$q=13$便得到$98\times13$。低阶系数压缩了谱包络，不能当成文字或音素标签。若编码器随后降采样4倍，在本章另行约定的向上取整输出规则下只有25个编码步；98个特征帧、25个编码步和若干文字符号是三个不同长度。

\subsubsection{声学特征增强：SpecAugment}
\label{sec:nlp-speech-specaugment}

识别器可能过分依赖某一频带或局部几帧。SpecAugment在训练输入的特征矩阵上遮蔽局部区域，让模型利用其余信息预测同一转写；原方案还包含时间扭曲。它改变的是声学观测，不是重新产生一份独立录音。

以零均值化后的$98\times80$特征为例，教学实现从整数集合$\{0,\ldots,F\}$均匀采样频带宽度$f$，再从$\{0,\ldots,80-f\}$采样起点$j_0$，令所有帧的第$j_0$至$j_0+f-1$列为0。时间掩蔽同样采样长度$r\leq R$和起点$i_0$，将第$i_0$至$i_0+r-1$行置零。若$f=8,j_0=20,r=10,i_0=30$，列20至27与行30至39被遮蔽；索引在这里均从0开始。置零是否代表平均特征，取决于之前的归一化，不能把零特征误称为真实静音。每轮训练重新采样掩蔽，损失仍使用原转写，参数由识别目标更新\scite{nlp-speech-park2019}

加噪在波形上叠加声音，混响模拟传播后的拖尾，变速改变时间尺度；它们与特征掩蔽作用在不同对象上。变速后文字通常可沿用，原时间戳却必须按变速比例变换；混入另一位清晰说话者时，是否仍只标主说话者需要明确协议。掩蔽过强若删除了日期的全部证据，训练标签虽未改动，样本已成为更难甚至含歧义的条件预测。增强强度由开发集选择，独立评价保留真实未见来源和声音条件。

\subsection{声音与文字的对应}
\label{sec:nlp-speech-correspondence}

转写通常没有说明每个字从哪一帧开始。设特征矩阵$\bm{X}$经编码器形成$T$步表示，目标文本为$y=(y_1,\ldots,y_U)$，符号来自词表$V$。声音的长短会随语速改变，$T$和$U$也不相等。下面的实现分别通过声学状态路径、路径求和或条件生成建立对应；流式能力还取决于编码器允许看见的声音范围，并不由损失名称决定。

\subsubsection{声学状态与解码图：HMM路线}
\label{sec:nlp-speech-hmm}

一种处理变长对应的办法，是让每个发音单位在若干声学状态中停留，再按允许的顺序前进。\term{隐马尔可夫模型}（hidden Markov model, HMM）的状态序列承担这项工作，其基本推断见第\ref{sec:pgm-dynamic-discrete}节。识别应用还需把状态连接到音素、词和语言约束。HMM--GMM用高斯混合分布$p(\bm{x}_i\mid s)$评价特征在状态$s$下的似然；混合权重、均值、协方差和转移概率由配对转写所允许的隐路径估计，词典发音则来自人工词典或另行核查的发音模型。

用对数分数表示，维特比状态量满足
\begin{equation}
 v_i(s)=\log p(\bm{x}_i\mid s)+
       \max_{r\to s}\{v_{i-1}(r)+\log a_{r,s}\}.
 \label{eq:nlp-speech-hmm-viterbi}
\end{equation}
起始状态取对数概率0，其余取$-\infty$；保存取最大值的前驱，在允许终态中选最佳分数并回溯。字词边界的语言模型分数也加到路径上。教学词典若把“hi”映射成两个音素/h/、/ay/，可允许路径$h_1,h_1,h_2,ay_1,ay_2,ay_2$；回溯得到状态后，沿词典弧读出“hi”，不能把六个状态直接拼成六个字。

Kaldi的$H\circ C\circ L\circ G$解码图把这些约束组合起来：$H$由转移标识映射到上下文相关音素，$C$处理音素上下文，$L$把音素映射到词，$G$提供词序列约束。声学分数在线附加到对应转移，搜索最终读取词标签。这里的转移标识还编码声学分布索引等信息，不能简化成“一个HMM状态就是一个词”\scite{nlp-speech-kaldi}

HMM--DNN保留解码接口，用神经网络给状态后验$p(s\mid\bm{x})$；在常见混合接口中用$\log p(s\mid\bm{x})-\log p(s)$作声学似然分数，忽略对候选共有的项。状态监督可由已有系统强制对齐取得，也可继续用序列目标训练。发音词典没有“海棠楼”时，即使声学模型听到了相近声音，受限解码图也可能无法输出该词；扩词典与训练声学模型解决的是不同问题。搜索剪枝过窄造成的丢失，还应与模型本身评分错误分开。

\subsubsection{对齐路径求和：CTC}
\label{sec:nlp-speech-ctc}

如果只有声音与整句文字，不想先指定一个唯一帧级对齐，可以把所有合法对应的概率加起来。\term{连接时序分类}（connectionist temporal classification, CTC）正是这样定义文本概率。它在词表$V$外增加空白符$\varnothing$，每个编码步给出分布$q_t(k)=p(k\mid\bm{X})$。空白只表示此步没有输出新符号，不等于实际静音。对长度$T$的路径$\pi$，先合并相邻相同符号，再删除空白，记此变换为$\mathcal{B}$，于是
\begin{equation}
 p(y\mid\bm{X})=
 \sum_{\pi:\mathcal{B}(\pi)=y}\prod_{t=1}^{T}q_t(\pi_t),
 \qquad L_{\mathrm{CTC}}=-\log p(y\mid\bm{X}).
 \label{eq:nlp-speech-ctc-probability}
\end{equation}
乘积假定给定输入后各步的路径输出条件独立，不要求波形帧相互独立；双向编码器仍可让每个$q_t$依赖整段声音。训练通过目标概率反向传播到输出头与编码器，不需要逐帧人工标签\scite{nlp-speech-graves2006}

操作顺序很重要：路径$(a,a,\varnothing,a)$先变成$(a,\varnothing,a)$，再变成$aa$。若先删空白再折叠，就会错误地变成$a$。目标中相邻相同符号必须由至少一个空白隔开。令$r(y)=\sum_{u=2}^{U}\mathbf{1}[y_u=y_{u-1}]$，非空目标存在合法路径的条件是$T\geq U+r(y)$；这是编码器输出长度的约束，不是原始采样数约束。

直接枚举路径随$T$指数增长。给目标前后及符号之间插入空白，形成扩展标签
$z=(\varnothing,y_1,\varnothing,\ldots,y_U,\varnothing)$，
长度$s_{\max}=2U+1$。定义$\alpha_{t,s}$为到第$t$步结束、位于扩展标签位置$s$的所有合法部分路径概率之和。对$U>0$，初值为
\begin{equation}
 \alpha_{1,1}=q_1(\varnothing),\quad
 \alpha_{1,2}=q_1(y_1),\quad
 \alpha_{1,s}=0\quad(s>2).
 \label{eq:nlp-speech-ctc-initial}
\end{equation}
令越界状态的概率为0。到达$s$可以停留、从前一位置进入，或者跳过中间空白；最后一种仅在当前标签非空白且与前两位不同的时候允许：
\begin{equation}
\begin{aligned}
 \delta_s&=\mathbf{1}[s>2,\ z_s\ne\varnothing,\ z_s\ne z_{s-2}],\\
 \alpha_{t,s}
 &=q_t(z_s)\bigl(\alpha_{t-1,s}+\alpha_{t-1,s-1}
                    +\delta_s\alpha_{t-1,s-2}\bigr),\\
 p(y\mid\bm{X})&=\alpha_{T,2U}+\alpha_{T,2U+1}.
\end{aligned}
\label{eq:nlp-speech-ctc-forward}
\end{equation}
结尾既可停在最后一个标签，也可停在尾空白，不能把整行状态全部相加。跨两位限制防止把目标$aa$中的两个$a$走成一个持续重复。空目标只有全空白路径，概率为$\prod_tq_t(\varnothing)$；若$T<U+r(y)$，目标概率为0，不能把无穷损失当成可正常学习的样本。

表\ref{tab:nlp-speech-ctc-numbers}给出三步完整分布及目标$ab$的前向量。第三步到位置4时，前三个可进入的状态量为$0.04,0.16,0.45$，乘以$q_3(b)=0.1$得到$0.065$；再加尾空白$0.020$，目标概率为$0.085$。

\begin{table}[htbp]
  \centering
  \begin{tabular}{@{}crrrrrrrr@{}}
    \toprule
    & \multicolumn{3}{c}{逐步分布} & \multicolumn{5}{c}{$z=(\varnothing,a,\varnothing,b,\varnothing)$}\\
    $t$ & $\varnothing$ & $a$ & $b$ & $\alpha_{t,1}$ & $\alpha_{t,2}$ & $\alpha_{t,3}$ & $\alpha_{t,4}$ & $\alpha_{t,5}$\\
    \midrule
    1 & .5 & .4 & .1 & .5 & .4 & 0 & 0 & 0\\
    2 & .4 & .5 & .1 & .2 & .45 & .16 & .04 & 0\\
    3 & .5 & .4 & .1 & .1 & .26 & .305 & .065 & .020\\
    \bottomrule
  \end{tabular}
  \caption{CTC教学计算。逐步分布来自构造，前向量由递推得到；状态列不是归一化后的概率分布。}
  \label{tab:nlp-speech-ctc-numbers}
\end{table}

同一结果可以完全枚举核对：折叠为$ab$的五条路径是
$(a,a,b)$、$(a,b,b)$、$(a,\varnothing,b)$、$(\varnothing,a,b)$和$(a,b,\varnothing)$，
概率分别为$0.020,0.004,0.016,0.025,0.020$，总和$0.085$，负对数为$2.4651$。目标$aa$在这三步中仅有$(a,\varnothing,a)$，概率为$0.064$。这个对照同时检查了路径求和、重复分隔和结束状态。

实际计算使用对数域。令$\ell_{t,s}=\log\alpha_{t,s}$，则把式\eqref{eq:nlp-speech-ctc-forward}的加法改为
$\operatorname{LSE}(v_1,\ldots,v_k)=m+\log\sum_j\exp(v_j-m)$，
其中$m=\max_jv_j$，不可达项取$-\infty$，全不可达时直接返回$-\infty$。最后对两个终态作LSE并取负号。这样只需$O(TU)$个状态更新；只算概率可滚动保存两行，训练还需保留或重算求梯度所需状态。

解码时，每步取最大概率再折叠，得到的是最佳单条路径对应的文本，而不是最高总概率文本。构造两步相同分布$q(\varnothing)=0.40,q(a)=0.35,q(b)=0.25$。贪心选全空白，路径概率$0.16$；文本$a$却汇总$(a,a),(a,\varnothing),(\varnothing,a)$三条路径，概率$0.1225+0.14+0.14=0.4025$，高于空文本。不能在得到贪心文本后只重算这一个候选就宣称找到了最佳文本。

前缀束搜索将对应同一文字前缀的路径合并，但必须区分末步是空白还是非空白。令$b_t(l)$与$n_t(l)$分别为前$t$步折叠为$l$且以空白、非空白结束的概率；初始$b_0(\epsilon)=1$，其余为0。算法\ref{alg:nlp-speech-prefix}是按本节概率定义推导的教学实现，未剪枝时精确合并所有前缀，保留前$K$项后才成为近似搜索。

\begin{algorithm}[CTC前缀的空白与非空白合并]
  \label{alg:nlp-speech-prefix}
  \begin{algorithmic}[1]
    \Require 分布$q_1,\ldots,q_T$，词表$V$，束宽$K$
    \Ensure 按合并概率选择的文字序列
    \State 初始化$b_0(\epsilon)=1,\ n_0(\epsilon)=0$，束为$\{\epsilon\}$
    \For{$t=1,\ldots,T$}
      \State 新映射$b',n'$全部置0
      \For{旧束中的每个前缀$l$}
        \State $s\gets b_{t-1}(l)+n_{t-1}(l)$
        \State $b'(l)\mathrel{+}=s\,q_t(\varnothing)$
        \For{$c\in V$}
          \If{$l\ne\epsilon$且$c=\operatorname{last}(l)$}
            \State $n'(l)\mathrel{+}=n_{t-1}(l)q_t(c)$
            \State $n'(lc)\mathrel{+}=b_{t-1}(l)q_t(c)$
          \Else
            \State $n'(lc)\mathrel{+}=s\,q_t(c)$
          \EndIf
        \EndFor
      \EndFor
      \State 合并全部同名前缀后，按$b'(l)+n'(l)$保留前$K$项
      \State 将保留项赋给$b_t,n_t$；平分时按固定词表次序处理
    \EndFor
    \State \Return $\argmax_l\{b_T(l)+n_T(l)\}$
  \end{algorithmic}
\end{algorithm}

同一末字再次出现时，非空白结尾只能延续当前字，空白结尾才能形成新的重复字，这就是两类状态都不可省的原因。对上述两步例子，前缀$a$在第一步为$(b_1,n_1)=(0,0.35)$；第二步得到$(0.14,0.2625)$，合计$0.4025$。束宽、词表截断与外部语言约束都会改变近似结果，应与CTC本身的文本概率分开报告。

图\ref{fig:nlp-speech-alignment}用两条折叠为同一文本的路径说明：符号出现在哪一帧仍是不确定的。CTC目标只监督文字，并不要求其高概率尖峰等于人工标注的发音起止点。

\begin{figure*}[htbp]
  \centering
  % Adapted from the academic-drawing TikZ template; dimensions are final size.
\begin{tikzpicture}[
  x=1mm,y=1mm,
  every node/.style={font=\figurefont\small,align=center,text=BookInk},
  cell/.style={draw=BookRule,line width=.8pt,minimum width=13mm,
    minimum height=9mm,inner sep=1mm},
  label/.style={anchor=east},
  flow/.style={book arrow}
]
  \foreach \i in {1,...,6} {
    \node[cell,fill=FigureInput!12] at ({27+17*\i},0) {$\bm{h}_{\i}$};
  }
  \node[label] at (31,0) {编码步};
  \node[label] at (31,-15) {路径A};
  \node[label] at (31,-30) {路径B};
  \foreach \i/\v in {1/八,2/八,3/{$\varnothing$},4/八,5/{$\varnothing$},6/{$\varnothing$}} {
    \node[cell,fill=FigureModel!10] at ({27+17*\i},-15) {\v};
  }
  \foreach \i/\v in {1/{$\varnothing$},2/八,3/{$\varnothing$},4/八,5/八,6/{$\varnothing$}} {
    \node[cell,fill=FigureModel!10] at ({27+17*\i},-30) {\v};
  }
  \draw[flow] (137,-15) -- (145,-15);
  \draw[flow] (137,-30) -- (145,-30);
  \node[anchor=west,text=FigureOutput] at (146,-15) {八八};
  \node[anchor=west,text=FigureOutput] at (146,-30) {八八};
  \draw[BookRule,line width=.8pt] (5,-42) -- (160,-42);
  \node[label] at (31,-54) {独立标注};
  \node[anchor=west,align=left] at (37,-54)
    {第一个“八”：0.02--0.10\,s，说话者甲\\
     第二个“八”：0.13--0.22\,s，说话者甲};
\end{tikzpicture}

  \caption{同一声学序列的两条CTC教学路径均折叠为“八八”，但符号出现的位置不同。空白为$\varnothing$，相邻重复先合并，再删除空白。最下方的时间区间与说话者是另行给定的示意标注，不是由两条路径证明的精确时间真值。}
  \label{fig:nlp-speech-alignment}
\end{figure*}
\FloatBarrier

\subsubsection{注意力条件序列生成：LAS}
\label{sec:nlp-speech-las}

CTC每步不显式读取已输出文字。若希望后一个字符的概率依赖前面已说出的字符，可以把识别写成声音条件下的序列生成。Listen, Attend and Spell（LAS）由声音编码器listener和注意力解码器speller组成：前者把特征序列变成较短的表示序列，后者每次读取与当前输出有关的声音信息，再预测一个字符。原论文的金字塔双向LSTM逐层拼接相邻时间步，减少注意力需要读取的位置数\scite{nlp-speech-chan2015}

记编码结果为$\bm{h}_1,\ldots,\bm{h}_T$，解码状态为$\bm{s}_u$。解码器根据前一状态、前一字符和前一上下文更新$\bm{s}_u$，计算相关分数$e_{u,t}$及权重
$a_{u,t}=\exp e_{u,t}/\sum_j\exp e_{u,j}$，
再汇总$\bm{c}_u=\sum_ta_{u,t}\bm{h}_t$。输出头以$[\bm{s}_u;\bm{c}_u]$给字符分布，目标包含结束符：
\begin{equation}
 L_{\mathrm{LAS}}=-\sum_{u=1}^{U+1}
       \log p_\theta(y_u\mid\bm{X},y_{<u}),
 \qquad y_{U+1}=\texttt{EOS}.
 \label{eq:nlp-speech-las}
\end{equation}
配对录音与文字训练编码器、注意力和输出头；教师强制使用正确前缀，推理从开始符出发，必须使用自己已经生成的前缀。原文还实验了训练中采样前一输出的策略，以缓解这项差异。

仍用“八八”短录音，LAS的两个$a$类字符可以分别由两个输出步生成，不用插入CTC空白。假设候选“八／八／结束”的三个条件概率为$0.7,0.6,0.8$，该候选概率为$0.336$；另一候选必须按它自己的前缀重新计算，不能只替换一个孤立字符概率。束搜索保存多个前缀，遇到结束符移入完成集合，并设置最大长度防止重复生成。注意力权重是读取声学信息的内部权重，不直接是字级时间真值；原始双向编码与全段注意力读取了未来声音，离线转写准确也不说明能在用户说完前交付。

\subsubsection{多任务声文生成：Whisper}
\label{sec:nlp-speech-whisper}

同一段中文声音既可要求中文转写，也可要求英语译文；没有任务条件，这两个目标无法区分。Whisper把语言、任务和是否输出时间信息编码为解码器词元，与声音特征一起规定生成目标。其原始编码器读取约30秒窗口的log-Mel特征，文本解码器按已有前缀输出后续词元，结构使用第\ref{chap:transformer}章的编码器与解码器机制\scite{radford2023whisper}

训练样本由音频片段及配对文字构造，任务前缀标明转写或到英语的翻译；语言标识、无语音标识和时间标识也进入多任务格式。部分样本把前段转写作为上下文。损失仍为目标词元的条件负对数概率，但“正确输出”现在由任务条件共同决定。时间词元表示当前窗口内的相对位置，原格式按20\,ms量化；量化间隔只是可表达的网格，不是20\,ms的边界准确率。

推理时固定语言与任务条件，贪心或束搜索生成文字和所需特殊词元，解析出片段并判断结束。长录音由窗口推进处理，已解码片段可用于提示下一窗口，窗口局部时间须加上其在录音中的偏移。“读书会改至四月十八日”在转写任务下应保留中文，在翻译任务下可以产生“April 18”；若把任务前缀设错，输出语言改变不一定是声学识别错误。词级时间戳还需实现规定的对齐过程，不能把所有时间输出都归于原始时间词元。

静音、背景音乐和未听清片段应作为独立测试条件。若没有语音却生成“谢谢观看”，文字流畅不能提供声音依据；前段上下文还可能使错误在多个窗口间重复。应检查无语音判断、窗口边界、重复片段和新增内容，必要时拒绝无证据转写。固定窗口上的离线生成没有自动解决稳定提交和低延迟，下一种识别概率模型与后面的流式组织分别处理不同问题。

\subsubsection{历史条件与路径求和：RNN-T}
\label{sec:nlp-speech-rnnt}

既想保留输出历史，又想对声文对应求和，可以让每个预测同时读取声音位置与输出前缀。\term{循环神经网络转导器}（recurrent neural network transducer, RNN-T）在二维格点$(t,u)$上定义符号分布：$t$表示当前编码步，$u$表示已经输出的目标符号数。声学网络得到$\bm{h}_t$，预测网络根据开始标识和$y_{1:u}$得到$\bm{g}_u$，联合网络归一化得到$q_{t,u}(k)$。原论文把两个网络的输出分数相加再做softmax，后续联合网络可以采用更一般的可学习映射\scite{nlp-speech-graves2012}

这里采用$t=1,\ldots,T$、$u=0,\ldots,U$的索引。空白边从$(t,u)$到$(t+1,u)$，消耗一个声音编码步；非空白$y_{u+1}$从$(t,u)$到$(t,u+1)$，增加一个输出而不推进声音。一个声音步可以连续发出多个标签。路径恢复只删除空白，不合并重复标签，因而目标$aa$可以在同一时间步连续发出两个$a$；CTC的$T\geq U+r(y)$不适用于RNN-T。

固定目标后，$\alpha_{t,u}$汇总进入格点前已产生$u$个目标符号的路径概率。定义$b_{t,u}=q_{t,u}(\varnothing)$、$e_{t,u}=q_{t,u}(y_{u+1})$，初始$\alpha_{1,0}=1$，越界项为0。递推与结束概率为
\begin{equation}
\begin{aligned}
 \alpha_{t,u}
 &=\alpha_{t-1,u}b_{t-1,u}
   +\alpha_{t,u-1}q_{t,u-1}(y_u),\\
 p(y\mid\bm{X})&=\alpha_{T,U}b_{T,U}.
\end{aligned}
\label{eq:nlp-speech-rnnt-forward}
\end{equation}
在首列只允许空白进入，首行只允许标签进入；$(1,0)$使用初值，不重复套递推。最后还必须从$(T,U)$发出一次空白，才能消耗最后一个声音步。结束前已经输出全部目标时，仍可沿空白边向后推进，但不能继续生成额外目标。

反向量$\beta_{t,u}$汇总从格点到终点的剩余概率，初值$\beta_{T,U}=b_{T,U}$。逆序递推为
\begin{equation}
 \beta_{t,u}=b_{t,u}\beta_{t+1,u}
       +e_{t,u}\beta_{t,u+1},
 \label{eq:nlp-speech-rnnt-backward}
\end{equation}
其中只保留合法边，$u=U$没有目标标签边，$t=T$没有通往下一普通编码步的边；末尾空白已包含在初值。于是$\beta_{1,0}=p(y\mid\bm{X})$。例如标签边的后验使用
$\alpha_{t,u}e_{t,u}\beta_{t,u+1}/p(y\mid\bm{X})$，
空白边同理。前向与反向把目标负对数概率的梯度分配给各边，再更新声学、预测和联合网络；核心格点计算为$O(TU)$，实际实现还要考虑词表归一化成本。

构造$T=2,y=aa$的例子，词表只有$a$和空白。第一行三个格点的$(b,q(a))$依次为$(.4,.6),(.5,.5),(.7,.3)$，第二行依次为$(.3,.7),(.2,.8),(.9,.1)$。图\ref{fig:nlp-speech-rnnt-lattice}给出合法目标边。三条完整路径
$(a,a,\varnothing,\varnothing)$、
$(a,\varnothing,a,\varnothing)$、
$(\varnothing,a,a,\varnothing)$
的概率为$0.189,0.216,0.2016$，总和$0.6066$。递推也得到$\alpha_{1,:}=(1,.6,.3)$、$\alpha_{2,:}=(.4,.58,.674)$，末尾乘$.9$得到同一结果。CTC在两步内不能产生$aa$，此例明确显示了两种路径空间的区别。

\begin{figure}[htbp]
  \centering
  % Target-specific probability lattice, not a measured alignment.
\begin{tikzpicture}[
  x=1mm,y=1mm,
  every node/.style={font=\figurefont\small,text=BookInk},
  state/.style={circle,draw=FigureModel,fill=FigureModel!10,
    minimum size=5mm,inner sep=0pt,line width=.9pt},
  time/.style={book arrow,draw=BookBlue},
  token/.style={book arrow,draw=BookTeal},
  edge label/.style={fill=white,inner sep=1pt}
]
  \foreach \u in {0,1,2} {
    \node[state] (a\u) at (20,{19*\u}) {};
    \node[state] (b\u) at (62,{19*\u}) {};
    \node[anchor=east] at (9,{19*\u}) {$u=\u$};
  }
  \node at (20,-10) {$t=1$};
  \node at (62,-10) {$t=2$};
  \draw[time] (a0) -- node[edge label,above] {$\varnothing:\ .4$} (b0);
  \draw[time] (a1) -- node[edge label,above] {$\varnothing:\ .5$} (b1);
  \draw[time] (a2) -- node[edge label,above] {$\varnothing:\ .7$} (b2);
  \draw[token] (a0) -- node[edge label,right] {$a:\ .6$} (a1);
  \draw[token] (a1) -- node[edge label,right] {$a:\ .5$} (a2);
  \draw[token] (b0) -- node[edge label,right] {$a:\ .7$} (b1);
  \draw[token] (b1) -- node[edge label,right] {$a:\ .8$} (b2);
  \node[draw=FigureOutput,fill=FigureOutput!12,minimum height=8mm,
    inner sep=2mm] (end) at (97,38) {终点};
  \draw[time] (b2) -- node[edge label,above] {$\varnothing:\ .9$} (end);
  \node[anchor=west] at (5,-22) {起点$(1,0)$概率为1；终点概率为0.6066};
\end{tikzpicture}

  \caption{两编码步、目标“$aa$”的RNN-T格点。横向空白推进时间，纵向$a$推进输出；标注为教学边概率。只有从右上角发出的末尾空白进入终点，重复标签没有折叠。}
  \label{fig:nlp-speech-rnnt-lattice}
\end{figure}
\FloatBarrier

贪心解码在当前格点取最大概率符号：空白使$t$加1，非空白追加到前缀、更新预测网络状态并保持$t$。束搜索则扩展多个前缀，将相同时间和相同前缀的等价路径概率相加，缓存对应预测网络状态。每步最多发多少符号、全句最大输出长度是避免搜索无界扩展的实现限制，不是原概率模型的CTC式长度约束；达到上限后强制推进也会改变搜索结果。

原始RNN-T使用双向声学编码器，并非天然在线。流式使用必须限制编码器看见的未来，例如因果计算或规定块长与前瞻，并保存跨块缓存。短停顿后继续说“十八日”时，编码器可能已积累声音却暂不发字，或在同一步发出多个字；符号位置、稳定提交和墙钟等待仍需分别记录。

\subsubsection{稳定前缀提交：Whisper-Streaming}
\label{sec:nlp-speech-whisper-streaming}

整段识别器也可通过缓冲和反复识别接入实时声音，但必须决定哪些结果还能改。Whisper-Streaming采用局部一致性（LocalAgreement）：比较连续几次、不断增长的声音输入对应的识别结果，只提交稳定的公共前缀。其LocalAgreement-2实现比较相邻两次的未确认文字；尾部保留在缓冲中等待更多证据。这是外部流式组织，不是原始Whisper模型自带的概率保证\scite{nlp-speech-machacek2023}

每轮收集至少规定长度的新声音后，识别当前声音缓冲；跳过已确认的文字部分，对上次和本次剩余序列求最长公共前缀，更新提交边界。已确认部分可能仍留在声音缓冲内提供上下文，但其新识别变化不作为对外撤回。待确认句界等条件满足后才裁掉旧声音，并用已提交文字保持跨句上下文。若计算本身慢于最低块长，下一轮处理积累的声音，所以实际等待由输入与计算共同决定。

设三次教学假设依次为“改到／四月／十七日”、“改到／四月／十八日”、“改到／四月／十八日／下午”。第一次尚无可比较结果；第二次提交“改到／四月”，保留日期；第三次才提交“十八日”。若把第一次末词直接播放成语音，后来修改屏幕文字也不能撤销听者已经听到的日期。若连续两次都误识为“十七日”，局部一致性照样会提交错误，它判断的是稳定性而不是真实性。

评价时固定比较单位、标点规范、缓冲裁剪和块长，分别统计最终字错误率、暂定文字修订量与单位提交延迟。删除的参考字没有提交时刻，应同时报告遗漏，不能通过只平均已识别字的延迟隐去漏识。固定块策略便于限定调度，局部一致性则随假设变化等待；两者都必须在真实计算预算下验证。

\subsection{解码、分段与转写整理}
\label{sec:nlp-speech-postprocessing}

取得字符后，系统还需按任务契约输出词、句子、时间和说话者。文字错误率以编辑距离为基础：设参考长$n$、假设长$m$，$d_{i,j}$表示把参考前$i$项变成假设前$j$项所需的最少编辑数，
\begin{equation}
 d_{i,j}=\min\{d_{i-1,j}+1,\ d_{i,j-1}+1,\
              d_{i-1,j-1}+\mathbf{1}[y_i\ne\hat y_j]\},
 \quad d_{i,0}=i,\ d_{0,j}=j.
 \label{eq:nlp-speech-edit-distance}
\end{equation}
回溯最短路径得到替换$s$、删除$d$、插入$i$，词错误率$\mathrm{WER}=(s+d+i)/n$按词计数，字错误率$\mathrm{CER}$按字计数。分词、数字规范和标点是否参与计分必须固定；空参考要单独报告插入数，不能除以0。

参考“四月十八日”的5个字，假设“四月十七日”，只有1次替换，CER为20\%；整句很长时同一个日期错误可能只占很小比例，却足以使业务失败。时间边界用与人工标注的差值或约定容差评价；说话者分段需在最佳标签置换后统计漏检、误检和混淆，并声明静音、交叠与边界宽容区的口径。这些误差不会自动进入普通WER。

\subsubsection{语言模型浅层融合与候选重排}
\label{sec:nlp-speech-fusion}

声音证据含混时，独立文本语料可以帮助比较“海棠楼”和相近发音候选。\term{浅层融合}（shallow fusion）在解码搜索中相加识别与语言模型分数，例如
\begin{equation}
 S(y)=\log p_{\mathrm{asr}}(y\mid\bm{X})
       +\lambda\log p_{\mathrm{lm}}(y)+\gamma |y|.
 \label{eq:nlp-speech-fusion}
\end{equation}
语言模型从允许的文本语料估计，$\lambda$与插入奖励$\gamma$在独立开发集上按任务错误选定，不由测试答案调节。CTC项先按概率求和合并同文本路径，再加入一次该文本的语言模型对数分数；前缀实现仅在新文字单位产生时更新相应语言模型增量。由于同一文本的语言模型因子可从路径概率求和中提出，等价的路径加权实现也成立。SpecAugment论文的识别实验也给出了独立语言模型融合及开发集选权重的实现实例。

教学候选甲、乙的识别对数分数为$-2.0,-2.4$，语言模型对数分数为$-3.0,-1.0$，长度相同。取$\lambda=.3$时总分分别为$-2.9,-2.7$，乙胜出；这只是改变选择，不证明乙确实在录音中。专业词表可用受控加分或词典约束接入，但也可能把别的发音吸附到常用专名。

候选重排先用识别器形成固定的$n$个完整候选，再重算同类组合分数，只能在已有集合中选择。若“海棠楼”在初次搜索中已被剪掉，后续语言模型再强也找不回来。比较时应固定候选预算，并检查错误来自候选未形成，还是形成后选错；识别分数、语言分数和长度的尺度必须一起校准。

\subsubsection{给定文本的CTC与HMM强制对齐}
\label{sec:nlp-speech-forced-alignment}

\term{强制对齐}（forced alignment）固定文本，只寻找它在声音中的对应位置。CTC可复用扩展标签图，将式\eqref{eq:nlp-speech-ctc-forward}中的求和换成最大值、记录前驱，并在两个合法终态中回溯最佳路径。HMM则把给定词序列展开成发音和状态图，在其允许转移中执行维特比搜索。对齐器参数来自已训练识别模型，并非每次对齐都重新估计声音模型。

假设编码步间隔为40\,ms，起点偏移为0，回溯路径是$(\varnothing,a,a,\varnothing,b,b)$。若约定以占用网格单元的外边界表示区间，可报告$a:[.04,.12)$、$b:[.16,.24)$秒。另一实现若用窗口中心或把空白分配给邻字，会得到不同边界；特征窗、卷积感受野和降采样偏移也要纳入换算。CTC尖峰占用的单元尤其不能直接当成整个音的持续时间。

目标含重复字符时必须保留空白分隔；停顿可吸收到空白或HMM静音状态。给错文字仍可能找到低分的近似路径，因而“强制对齐成功”不是转写正确的证据。应用中可据归一化分数和异常长度筛查，再对照人工边界与原声音核验，不把模型时间戳当成新的监督真值。

\subsubsection{说话者嵌入与聚类：x-vector}
\label{sec:nlp-speech-xvector}

“谁在何时说话”先需要把活动声音分成适当片段，再比较说话者。\term{说话者嵌入}是对一个声音片段的固定长度表示，希望同一说话者的片段更接近。x-vector用说话者身份监督训练帧级网络，将变长帧表示的均值和标准差汇总为统计池化向量，再通过片段级网络预测训练说话者；使用时取内部片段表示而非身份softmax输出。它原本服务说话者识别，下面与聚类组成分段的流程是教学组合\scite{nlp-speech-snyder2018}

语音活动检测取得候选片段后，提取并归一化嵌入$\bm{e}_i$，用余弦相似度$\bm{e}_i^\top\bm{e}_j$或另行学习的评分模型比较。凝聚层次聚类从一片一组开始，合并相似度最高的两组，按开发集阈值或已知组数停止。假设三片段的相似度为$s_{1,2}=.91,s_{1,3}=.18,s_{2,3}=.22$，取组间平均相似度和阈值$.7$，先合并1、2，合并后与3的相似度$.20$，因此保留两组。再把片段组映射回其起止时刻，得到匿名说话者A、B。

训练身份分类只提供可比较的表示，不保证新说话者必然分开。短片段统计不稳，两人交叠时单一嵌入可能混合身份；需要允许重叠标签或增加分离、重分段机制。聚类A不能未经核验就命名为某位真实用户。评分时匿名标签须先匹配参考身份，并分别报告分段边界和身份混淆，不能因组数正确而宣称分段正确。

\subsubsection{标点恢复与反向文本规范化}
\label{sec:nlp-speech-itn}

原始识别可能输出“明天可以吗不可以后天”。标点恢复给字符之后的边界预测逗号、句号或问号，也可用条件生成形成整理文本。序列标注训练使用去除标点的输入和原标点位置标签；生成训练使用口头转写与整理文本对，逐词元计算交叉熵。推理恢复标点后，还要检查是否把“不可以”错误归到前一句。离线模型若看见后文，不能把其结果与只见当前前缀的在线模型直接比较。

\term{反向文本规范化}（inverse text normalization, ITN）把读出的形式恢复成规定书面形式，如“四月十八日下午两点”恢复为“4月18日14:00”。可先识别日期、时间、数量和号码类型，再用规则转写；也可训练类型标签与条件生成模型。类型决定读法的逆映射：“二零一”可能是房间201或一串编号，不能无条件按普通数量处理；“四月十六日十八点前”中的“前”必须保留。

ITN通常不是一一可逆：相同口头形式可能对应不同大小写、简称或格式，需要上下文及规定输出风格。后处理只改变允许的书面形式，不应凭语言流畅性把听成的“十七日”改为“十八日”。评价既保留统一规范下的识别错误，也单独计整理后的数量、单位、句界与可读性；否则规范化可能掩盖内容差异。

\section{口语理解}
\label{sec:nlp-speech-understanding}

用户希望系统理解“后天三个人”，不一定需要看到完整转写。\term{口语理解}（spoken language understanding, SLU）把声音及允许的历史转为业务语义或交流作用。先识别再理解可以检查中间文字，直接声音理解则让声音表示直接参与语义预测。两条路线都需要约定语义目标；删除外部转写接口不等于删除声学歧义。

\subsection{意图、槽位与口语事实}
\label{sec:nlp-speech-semantics}

另设一个与读书会无关的预约教学请求：“明天两个人，改成后天三个人。”咨询发生于2025年4月17日，当前有效日期应为4月19日、人数为3；4月18日与人数2是被撤回值，不能同时留在有效记录里。意图说明请求做什么，槽位记录其参数，口语事实还可按第\ref{chap:nlp-information-extraction}章的来源与事件表示保存证据。

SLURP将声音、转写和场景、动作、实体等语义标注配对，并分析实体内容错误，使流水线中的识别与语义错误能够分别检查。它为评价组织提供依据，不表示该语料提出了下面所有分类和结构生成算法；完整来源见文献说明。

\subsubsection{意图分类与CRF槽位标注}
\label{sec:nlp-speech-intent-crf}

把转写字符或词元交给文本编码器，句级表示$\bm{h}_{\mathrm{sent}}$接softmax预测意图，逐位置表示$\bm{h}_i$给BIO标签打分。这里B表示槽位开始，I表示继续，O表示非槽位；日期、人数可分别设标签。独立头分别训练，也可共享编码器，以意图交叉熵和CRF负对数似然之和更新参数。CRF的标签转移与维特比推断见第\ref{chap:nlp-text-analysis}章。

设意图为$c$、槽位标签为$z$，联合教学目标可写
$L=-\log p(c\mid x)-\lambda\log p_{\mathrm{crf}}(z\mid x)$，
其中$x$是已取得的转写，$\lambda$由开发集选定。预测时在合法BIO路径中解码，再把跨度映射回原文字。用词元“明天／两／个／人／改成／后天／三／个／人”，可恢复两个日期跨度和两个人数跨度；接着利用“改成”的更正关系，把后一组值设为当前有效值，前一组标为撤回。日期解析使用咨询时点，不使用训练日期或通知发布日期。

这个恢复过程不能止于槽位F1。标签序列即使完全合法，也可能把“后天”规范为4月18日，或没有撤销“两个人”。应分别核查意图、跨度、规范值、更正关系和最终完整记录。将人工正确转写与识别一最佳分别送入同一语义模型，可以定位识别误差的影响；识别候选集若包含正确日期，还需检查语义重排是否选中，不因候选里有答案就计为成功。

\subsubsection{直接声音分类与语义记录生成}
\label{sec:nlp-speech-direct-slu}

若目标只是判断是否在预约，声学编码序列可以池化为$\bm{h}_{\mathrm{audio}}$，再用声音—意图配对样本训练分类头。均值池化以全部有效帧求平均；注意力池化则学习每帧权重并归一化。样本掩码须排除补齐帧，否则长短录音的无声填充也可能影响表示。输出类别由softmax最大值或预先校准的拒识阈值恢复。

日期和人数还需要更细的输出。可用声音编码器接结构化解码器，以目标记录$z$的词元负对数概率
$-\sum_j\log p(z_j\mid\bm{X},z_{<j})$训练，模式约束只允许合法字段名、类型和结束位置。推理将生成序列解析为记录，再检查必填字段、日期范围、人数类型和更正状态。上述请求的目标可以是
\begin{quote}
意图：修改预约；日期：2025-04-19；人数：3；撤回：日期2025-04-18、人数2。
\end{quote}
这些值的监督来自核验过的语义标注；合成声音可以补充训练，但需与真实声音来源分开记录。

若还要求证据定位，必须说明监督对应声学帧区间还是转写字符区间，不能把没有转写的输出称为文字跨度。直接生成“4月19日”不证明模型听清了“后天”，仍需对照声音和咨询时点检查。与流水线比较时固定声音、语义模式及训练资源，报告关键槽位和完整请求正确率；池化分类不提供槽位生成能力，记录生成也不自动获得更正关系的可靠性。

\subsection{话语行为与情绪表达}
\label{sec:nlp-speech-pragmatics}

“可以”既能回答可行性询问，也能授予许可或表示勉强接受。\term{话语行为}描述一句话在当前交流中的作用；情绪标签则按协议描述表达出的情绪类别或唤醒度、效价量表，不等于对真实心理状态的直接测量。输入可以包括声音、文字、历史、说话者和回应对象，各实验必须固定其可见范围。

标注协议应允许不确定或多标签，并记录标注者一致性。相同文字配不同语调能检查声音线索，相同声音配不同历史能检查语境；划分时按会话或说话者留出，防止模型记住某人的声音。类别不均衡时报告宏平均F1或各类召回的非加权平均，回归任务另报告指定量表下的误差与相关性，不能直接比较不同标签体系的一个总分。

\subsubsection{历史条件下的话语行为分类}
\label{sec:nlp-speech-dialogue-act}

一个可复现的分类输入是“前轮提问、当前文字、当前说话者、回应对象、已确认状态”。编码这些内容后，用交叉熵学习请求、告知、确认、否定、更正等标签；若一句话同时承担确认与补充告知，可改用多个二元标签及开发集阈值。标签来源是按同一协议标注的对话轮，而不是从分类器自己的输出反复生成所谓真值。

若前轮是“周五下午两点方便吗”，当前“可以”可标为确认时间；前轮是“我可以修改人数吗”，同样的回答可能表示许可。独立当前句分类器看到完全相同的输入，无法区分这两种条件；加入回应对象和前轮内容，分类的依据才充分。推理输出标签后，状态更新还需核对确认的是哪个值，不能把一句“可以”解释成对所有未决动作的授权。

历史截断若删掉前轮问题，系统可以输出不确定或询问“你是确认时间，还是允许修改人数”。评价将标签正确与对话后果分开：行为分类为“确认”仍可能确认了错误的日期。声音语气可作为附加条件，但不能代替缺失的回应对象。

\subsubsection{韵律特征、SVM与声文融合}
\label{sec:nlp-speech-prosody-classification}

同一句“可以”读得平缓或急促，声音中的节奏信息可能帮助解释表达。\term{韵律}包含音高变化、能量、时长和停顿等超越单个音素的组织。一个基线先从有效语音帧提取基频$F_0$、对数能量和频谱或倒谱特征，再统计均值、标准差、分位数、斜率、总时长与停顿比例，形成固定向量$\bm{v}$。无声或不规则振动段的$F_0$应标为缺失，并另留有声帧比例，不能把“无基频”直接当作0\,Hz参与均值。

以二分类标签$z\in\{-1,1\}$为例，标准化参数只在训练集估计，支持向量机（support vector machine, SVM）最小化
$\frac12\lVert\bm{w}\rVert_2^2+C\sum_i\max(0,1-z_i(\bm{w}^{\top}\bm{v}_i+b))$。
$C$由开发集选定，推理使用决策分数和阈值；逻辑回归则用相同特征最小化交叉熵。类别多于两个时规定多类方案，唤醒度或效价任务则改为回归损失。声学编码器池化可替代手工统计，声文融合进一步拼接声音向量与文本向量，或让文字位置读取声音序列后分类。

教学对照固定“可以”的文字，仅把时长从0.8\,s改为0.4\,s并改变音高走势，用来说明分类器看到的输入变化，不预设短声必然表示不耐烦。再固定声音，把背景从“确认参会”改成“被迫让步”，观察标签协议是否允许不同解释。只有在留出的声音与会话上比较纯文字、纯声音及融合模型，才能判断声音提供了什么增益；即使识别出某种情绪表达，恰当回应仍由对话目标和来源条件决定。

\section{语音合成}
\label{sec:nlp-speech-tts}

\term{语音合成}（text-to-speech, TTS）把文字及规定的声音条件转换成可播放波形。文字没有唯一声音实现：同一通知可以快读、慢读、强调新日期，也可在地点前停顿。但这些自由不包括任意改变日期、数量和否定。合成流程因此既要形成声学细节，也要保持所交付的语言内容。

图\ref{fig:nlp-speech-synthesis}把读法、发音、韵律、声学表示与波形分开，是为了定位错误，并不要求每个系统有七个独立模型。下面先决定“读什么”，再决定“怎样发声”，在波形形成后完成一次贯通核验。

\begin{figure*}[htbp]
  \centering
  % Vertical role flow; dashed boundary denotes roles that a model may combine.
\begin{tikzpicture}[
  x=1mm,y=1mm,
  every node/.style={font=\figurefont\small,align=center,text=BookInk},
  block/.style={book node,text width=51mm,minimum height=10mm},
  check/.style={text width=63mm,align=left,anchor=west},
  flow/.style={book arrow},
  check edge/.style={book arrow,draw=BookMuted}
]
  \node[block,draw=FigureInput,fill=FigureInput!12] (text)
    at (38,0) {书面文字\\重庆银行，10.5万元};
  \node[block,draw=FigureModel,fill=FigureModel!10] (read)
    at (38,-18) {规范读法\\十点五万元};
  \node[block,draw=FigureModel,fill=FigureModel!10] (phone)
    at (38,-36) {发音单位\\重：chong2；行：hang2};
  \node[block,draw=FigureModel,fill=FigureModel!10] (duration)
    at (38,-54) {时长与韵律\\语速、停顿、重音};
  \node[block,draw=FigureModel,fill=FigureModel!10] (mel)
    at (38,-72) {声学序列\\梅尔频谱};
  \node[block,draw=FigureOutput,fill=FigureOutput!12] (wave)
    at (38,-90) {波形\\采样率、采样数};
  \node[block,draw=FigureOutput,fill=FigureOutput!12] (listen)
    at (38,-108) {听觉核验\\内容与自然度分别检查};
  \foreach \a/\b in {text/read,read/phone,phone/duration,duration/mel,mel/wave,wave/listen} {
    \draw[flow] (\a.south) -- (\b.north);
  }
  \node[check] (value) at (90,-18) {数值不变：105000元\\不误读成10.5元};
  \node[check] (name) at (90,-36) {上下文决定多音字\\专名另核读音};
  \node[check] (length) at (90,-63) {时长扩展后再检查\\遗漏、重复与停顿};
  \node[check] (audit) at (90,-108) {回转识别只作筛查\\盲听、配对比较另行记录};
  \draw[check edge] (read.east) -- (value.west);
  \draw[check edge] (phone.east) -- (name.west);
  \draw[BookMuted,dashed,line width=.8pt] (8,-45) rectangle (70,-81);
  \draw[check edge] (70,-63) -- (length.west);
  \draw[check edge] (listen.east) -- (audit.west);
  \node[anchor=west] at (90,-85) {虚框：可由同一模型形成};
\end{tikzpicture}

  \caption{文本到声音的职责与检查。例句“重庆银行，10.5万元”只用于教学，不属于读书会通知。右侧检查分别定位金额、读音、时长和最终听觉内容；虚框内职责可以由一个声学模型联合实现，职责合并不消除相应错误。}
  \label{fig:nlp-speech-synthesis}
\end{figure*}
\FloatBarrier

\subsection{文本规范化与发音}
\label{sec:nlp-speech-normalization}

“10月5日”“10.5万元”“A区”“重庆银行”表面上都是字符，却需要不同的读出方式。文本规范化将书面形式变成可读形式，发音映射再指定音素或带声调的拼音等单位；前者应保持值和单位，后者应在上下文中选定读音。直接字符输入的声音模型也承担这些职责，只是错误更难在独立接口上观察。

\subsubsection{规则词典与WFST规范化}
\label{sec:nlp-speech-wfst-normalization}

一个规则实现先标出日期、金额、编号、缩写和普通文字，再按类型展开读法。带权有限状态转导器（weighted finite-state transducer, WFST）把这种过程写成带输入、输出和代价的有向路径：输入弧消费书面符号，输出弧产生读法单位，路径代价比较多个候选。规则及词典来自明确编写的语言约定，候选权重可人工设定或用已标注的书面—读出形式对估计，不必假定所有规则都由训练学得。

对“10月5日”，日期规则消费月份10和日期5，输出“十月五日”；对“10.5万元”，金额规则输出“十点五万元”，其规范值仍为105000元。输入“201”若被类型标为房间编号，可按项目读法规定输出“二零一”；若类型为计数，则可能输出“二百零一”。候选路径必须同时满足类型与上下文约束，再用最短路径取出读法；同一数字字符串不能仅凭其局部形状决定类型。

电话号码按位读、金额按数值读，正好说明规则冲突在哪里发生。对“A区”可规定按英文字母名称读A再读“区”，但未经说明不能扩成某个组织缩写的全称。词典缺项应保留可检查的未知或转交发音模型；规则有重叠时用优先级或权重消歧，并测试极端日期、小数、负号和单位组合。读得连贯却漏掉“万”，仍是内容失败。

\subsubsection{句级神经文本规范化}
\label{sec:nlp-speech-neural-normalization}

规则需要显式枚举写法，句子上下文还可能改变某个缩写的含义。句级神经规范化把整句书面文字作为输入，把读出形式作为目标，学习条件分布。Mansfield等比较窗口级与句级序列生成，并使用子词单位以及大小写、词性和位置等特征；原文实验不能据此当成中文多音字的验证\scite{nlp-speech-mansfield2019}

训练样本$(x,r)$可以来自人工标注，也可以由规则系统生成后筛查；后一种数据会继承规则错误。模型按教师强制最小化$-\sum_j\log p(r_j\mid x,r_{<j})$，推理以贪心或束搜索形成候选，拼回读出单位，遇结束符停止。与逐窗口规范化相比，整句模型可利用更远上下文，但也可能改变本来不需修改的文字。子词缓解固定输出词表对新组合的限制，不能保证未见金额一定正确。

对“预算10.5万元，10月5日确认”，合格候选应同时保持金额105000元与日期10月5日。可以把输入和候选中的类型、数值、单位解析成结构记录，要求金额与日期分别相等；解析失败的候选拒绝自动接受。这个确定性筛查能发现漏“万”，却不能发现所有听觉发音错误。正向规范化与第\ref{sec:nlp-speech-itn}节ITN的输入输出相反，训练对的方向不同，也不保证彼此完全可逆。

\subsubsection{字形到发音：G2P}
\label{sec:nlp-speech-g2p}

得到“重庆银行”后，系统还须决定“重”和“行”怎么读。\term{字形到发音}（grapheme-to-phoneme, G2P）把文字映射成发音单位序列。词典覆盖的词可直接取得候选读音，上下文分类器给多音候选评分；未登录词可由字符或子词编码器生成音素序列，以词典和人工发音标注训练条件交叉熵。推理在允许音素词表上搜索，随后核对字词与音素数量、音节边界和语种标识。

中文实现需约定分词、带调拼音或声韵母单位，以及变调和轻声是在前端标注还是由声音模型实现。这里的词典教学表示采用本调：“重庆”为\code{chong2 qing4}，“银行”为\code{yin2 hang2}，不把每个字的最常见读音当成唯一读音。通知里的“海棠楼”可核为\code{hai3 tang2 lou2}，“201”先经编号规范化再转为发音单位；\code{ML-0418}若需要朗读，必须事先规定字母和数字逐位读法。

直接字符TTS可把读音选择隐含在模型中，显式音素TTS则把这一选择暴露出来，便于修正专名。两者都可能因上下文不足选错发音。G2P解决音素身份，不独自决定停顿、句重音和整体语速；把正确音素用不合适的节奏说出，仍可能难以理解。

\subsection{时长、韵律与声学序列}
\label{sec:nlp-speech-acoustics}

文本或音素只有几十个单位，输出频谱往往有几百帧。声学模型必须决定各单位持续多久，并形成音高、能量和停顿。训练样本包括文本或发音、对应录音提取的梅尔频谱，以及允许使用的说话者或表达条件。频谱提取参数应与声码器一致，不能只因两者都有80个频带就直接拼接。

\subsubsection{注意力声文对齐：Tacotron 2}
\label{sec:nlp-speech-tacotron}

Tacotron 2从字符编码开始，每一步用注意力读取文字表示，以前一步声学帧和当前上下文预测下一梅尔帧。训练时有目标频谱，教师强制把正确上一帧送入解码器；推理只能反馈自己生成的帧。原模型采用位置敏感注意力，利用此前累计的注意力帮助向前读取文字，另由后处理网络预测频谱残差，并用停止标识判断生成是否结束。其完整来源列于本章文献说明。

设原始预测为$\widehat{\bm{M}}$、残差修正后为$\widehat{\bm{M}}^{+}$、目标为$\bm{M}$，帧预测损失包括两处均方误差，再加停止标识的二元交叉熵。损失同时训练字符编码器、注意力、解码器和后处理网络。原论文采用每步一帧，推理停止概率首次超过规定阈值时结束；应用还应设置最大帧数并标记异常终止，不能把截断当成正常完成。

对“四月十八日”可观察注意力是否依序经过各字，再检查输出中是否真的包含每个音节。注意力卡在“十”附近可能对应重复，提前到结束位置可能对应漏读“八日”；但单凭注意力图也不能判断可懂度，仍需听觉核验。与下一小节不同，Tacotron 2没有显式逐音素整数时长预测器。它的训练频谱监督也不能直接替代文本规范化：原文使用规范化后的英文文字，中文数字和专名需要本任务自己的读法条件。

\subsubsection{时长与韵律条件：FastSpeech 2}
\label{sec:nlp-speech-fastspeech}

若先给每个音素分配帧数，就可以在长度确定后并行预测频谱。FastSpeech 2把时长、音高和能量显式加入声音生成条件。训练时从配对录音取得这些量，推理时使用一同训练的预测器估计；它不再以教师模型合成频谱作为主要目标。原论文的声学输出仍是梅尔频谱，直接生成波形的是另行设计的FastSpeech 2s\scite{nlp-speech-ren2021}

训练准备先用强制对齐器将核验过的文字或音素对到录音，换算每个音素的整数梅尔帧时长$d_i$，使总和与目标频谱帧数一致；对齐失败需筛查，不能凭空分配。音高从波形估计，须区分有声与无声；能量可由STFT幅度向量的范数得到。原方法将连续音高曲线作连续小波变换，预测其多尺度表示，并恢复音高曲线及所需统计量；音高与能量量化后的嵌入加入隐藏序列。这里的时长、音高和能量均是从训练声音提取的监督，不是人工为每条文本指定的唯一正确表达。

设音素编码为$\bm{h}_1,\ldots,\bm{h}_n$，\term{长度调节器}按$d_i$重复第$i$个向量，使隐藏长度与频谱一致：
\begin{equation}
 \widetilde{\bm{h}}_t=\bm{h}_i
 \quad\text{当}\quad
 \sum_{j<i}d_j<t\leq\sum_{j\leq i}d_j,
 \qquad T_{\mathrm{mel}}=\sum_{i=1}^{n}d_i.
 \label{eq:nlp-speech-length-regulator}
\end{equation}
教学音素序列/p, a, t, a/的时长为$(2,3,1,2)$帧，则展开为
\[
 (\bm{h}_1,\bm{h}_1,\bm{h}_2,\bm{h}_2,
  \bm{h}_2,\bm{h}_3,\bm{h}_4,\bm{h}_4),
\]
共8帧。它只是检验长度扩展的短序列，不冒充自然语音中的典型音素时长。若帧移10\,ms，对应帧网格长80\,ms；加窗补边对最终波形长度的影响仍按声码器接口确定。

频谱重建误差、对数时长的均方误差、音高表示的预测误差与能量均方误差共同训练模型。为使零时长也有定义，一个明确的教学时长目标采用$\log(1+d_i)$，推理按
$\hat d_i=\max(0,\operatorname{round}(\exp\hat r_i-1))$
恢复整数；若任务要求每个非静音音素至少占一帧，还应显式设下界1，而不是悄悄丢掉音素。这一取整和下界是实现约定，不能与原论文所有代码变体混为一谈。训练长度调节器使用提取时长，频谱模型看到提取的音高、能量条件；推理则全部改用预测量，误差会互相影响。

语速控制可将预测时长除以速度因子后取整。对$(2,3,1,2)$取速度因子1.5、采用四舍五入且非静音下界1，得到$(1,2,1,1)$，共5帧；$8/1.5$不是整数，实际比例不会恰好等于1.5。停顿应作为明确的停顿单位增加，例如在第二、三音素之间插入2帧静音，而不是把邻音素改成另一个音。音高或能量调整可以改变重音线索，但音高变高不保证听者理解为正确重点。

扩展后必须检查每个应读音素是否保留、帧数是否合理、停顿是否落在允许位置，再把频谱交给声码器。对通知强调“十八日”时，允许改变时长与强度，不允许删掉“校外来宾”或“前”。与Tacotron 2相比，FastSpeech 2显式确定对齐长度并并行预测频谱，Tacotron 2在逐帧生成中隐式形成对齐；两者都可能发错音，也都需要最终内容核验。

\subsection{波形形成与内容核验}
\label{sec:nlp-speech-waveform}

梅尔频谱还不是可播放声音，它丢失了包括相位在内的细节。\term{声码器}（vocoder）在声学条件下恢复波形。一个可复现接口至少规定频谱提取、采样率、每帧对应采样数、补边与裁剪；波形时长由实际采样数除以采样率得到，不能把帧数直接当成秒数。

\subsubsection{条件逐采样生成：WaveNet}
\label{sec:nlp-speech-wavenet}

条件WaveNet把波形概率分成逐采样条件概率。设目标波形为$x_1,\ldots,x_N$，梅尔条件为$\bm{M}$，则
\begin{equation}
 p(x_{1:N}\mid\bm{M})
   =\prod_{n=1}^{N}p_\theta(x_n\mid x_{<n},\bm{M}),
 \qquad
 L_{\mathrm{wave}}=-\sum_{n=1}^{N}
       \log p_\theta(x_n\mid x_{<n},\bm{M}).
 \label{eq:nlp-speech-wavenet}
\end{equation}
扩张因果卷积读取已知采样，声学条件经上采样与波形时间分辨率对应。Tacotron 2中的修改版预测混合逻辑分布参数，用真实波形采样训练负对数似然；它还使用教师强制生成、与目标波形对齐的预测频谱训练声码器，使训练条件接近上游模型的频谱，而非默认所有条件都是真实频谱。

推理从规定的初始状态开始，逐个采样并反馈，直到产生接口要求的$N$个采样。教学设定采样率24\,kHz、帧移12.5\,ms，一帧网格对应300个采样，80帧对应24000个采样，即1秒；这里不另加窗尾，实际系统应说明裁剪规则。逐采样依赖使生成有$N$个顺序步骤，缓存能减少重复卷积，却不把自回归自动变成并行生成。

固定同一份频谱可以单独比较声码器引入的杂音和可懂度。但如果上游频谱本来对应“十七日”，形成自然波形也不会使它变成“十八日”。读法、对齐或频谱的内容错误应回到对应环节修正，不能只调整波形音质。

\subsubsection{并行波形生成：HiFi-GAN}
\label{sec:nlp-speech-hifigan}

HiFi-GAN用卷积生成器直接把梅尔序列上采样到波形长度，减少逐采样的顺序依赖。其训练样本是配对波形及从该波形提取的梅尔表示；转置卷积后接不同感受野的残差模块，生成器参数由声音判别与重建目标共同更新\scite{nlp-speech-kong2020}

多周期判别器把波形按若干周期重排，分别比较相同相位间隔上的模式；原方法采用2、3、5、7、11等周期。多尺度判别器同时观察原波形及平均池化后的不同尺度，检查连续局部模式和较长时间结构。令$\hat x=G(\bm{M})$、第$k$个判别器为$D_k$，生成器目标可概括为
\begin{equation}
\begin{aligned}
 L_G={}&\sum_k\mathbb{E}(D_k(\hat x)-1)^2
       +\lambda_{\mathrm{fm}}L_{\mathrm{fm}}
       +\lambda_{\mathrm{mel}}L_{\mathrm{mel}},\\
 L_{\mathrm{mel}}={}&\mathbb{E}\lVert
        \phi(x)-\phi(\hat x)\rVert_1 .
\end{aligned}
\label{eq:nlp-speech-hifigan}
\end{equation}
$\phi$为固定梅尔提取器，$L_{\mathrm{fm}}$为真实与生成波形在判别器各层特征上的归一化$L_1$距离。判别器自身最小化真实输出接近1、生成输出接近0的最小二乘损失，和生成器交替更新。对抗项检查可分辨的波形模式，特征匹配缓和中间表示差异，梅尔项约束声学内容；原HiFi-GAN的该重建项不能误写成多分辨率STFT损失。

推理只保留生成器，按已知上采样倍率产生波形并裁到约定长度。固定80帧频谱，比较WaveNet与HiFi-GAN时，应统一采样率、频谱定义、硬件、批量和计时边界，分别记录首段可用时间与整段生成时间。并行结构减少顺序步骤，不给任意硬件上的速度保证；判别器认为逼真，也不说明每个数字读对。

把流程贯通到共同通知，可以选取完整的待播文本：“机器学习读书会改至4月18日14:00，地点为海棠楼201。校外来宾须在4月16日18:00前提交报名信息。”规范化应保留活动时间与报名截止两个不同字段，产生“四月十八日下午两点”和“四月十六日十八点前”等允许读法；G2P核对“海棠楼”及编号，声学模型保留每个音节及适当断句，声码器按规定采样率恢复波形。

内容验收再按四项独立问题进行：听者是否听到4月18日14:00，是否听到海棠楼201，是否保留校外来宾这一对象，是否听到严格早于4月16日18:00的截止条件。回转ASR可以自动比较这些字段，但识别器可能把发错的音纠正成常用词，或者把正确声音误识，因此只能筛查。真人盲听转写、配对偏好与平均意见分（mean opinion score, MOS）应记录材料、评价者、量表和不确定性，内容正确、可懂度与自然度分别报告。

若只把已合成通知中的“十七日”改成“十八日”，应声明修改区间、需要保持的前后区间和允许的接续变化。新日期时长不同，就要重新对齐后续边界并检查拼接，不能把局部替换后的所有旧时间戳原样保留。声音词元模型可以把声学与波形职责合并，经编解码器恢复声音；音色相似、频谱接近或自然度高仍不替代上述内容核验。

\section{语音翻译}
\label{sec:nlp-speech-translation}

语音翻译把源语言声音转换成目标语言文字或声音，核心约束仍是第\ref{chap:nlp-text-generation}章讨论的意义保持。级联与直接描述实现路线，文字与声音描述交付对象，整段与流式描述可见信息和时间条件；这三个维度不能混成一组互斥模型类别。即使屏幕上的英文正确，目标声音读错时间仍是语音到语音任务失败。

\subsection{级联转换}
\label{sec:nlp-speech-cascade}

级联路线串起已经建立的接口：源语识别提供文字，文本翻译产生译文，需要时再合成为声音。CoVoST 2提供多语种声音与译文配对，可用来组织声音条件的翻译监督与评价；它不规定所有应用都必须沿级联运行，也不替代中文通知的专名和时间核验\scite{nlp-intro-wang2021covost}

\subsubsection{一最佳识别—翻译—合成}
\label{sec:nlp-speech-one-best}

设源声音为$\bm{X}$，源语文字为$x$，译文为$y$。一最佳级联依次搜索
\[
 \hat x=\argmax_x p_{\mathrm{asr}}(x\mid\bm{X}),
 \qquad
 \hat y=\argmax_y p_{\mathrm{mt}}(y\mid\hat x),
\]
再将$\hat y$及目标语言、声音条件交给TTS。两个搜索通常都使用有限候选近似。识别器用声音—转写监督，翻译器用平行文本，合成器用目标语言文字—声音配对；联合调优也不能省略接口的语言、分段、标点和专名约定。

以通知时间片段为例，正确源语转写“改至四月十八日下午两点”应译为“rescheduled to April 18 at 2 p.m.”等同义形式。若识别输出“四月十七日”，翻译器忠实译成“April 17”，局部翻译没有擅改意义，但端到端结果错了。把人工正确转写和识别转写分别送进同一个翻译器，可以测出识别误差带来的变化；对正确转写仍误译的部分，再定位翻译错误。TTS若把“2 p.m.”读成上午两点，还要单独归为目标声音内容错误。

中间记录应保存候选、版本和对应源声音区间。人工纠正“April 17”之前必须回听源声音或核对任务提供的依据，不能因为共同通知常见“April 18”就自动覆盖另一段实际谈论4月17日的声音。最终验收既核查译文的事实，也核查听者实际获得的目标声音，不能以生成了波形文件代替完成。

\subsubsection{识别多候选传递与重排}
\label{sec:nlp-speech-nbest-translation}

一最佳识别过早抛弃了其他解释。若声学上“十七”与“十八”都可能，可以保留识别$n$个候选，每个再翻译，按源语识别、条件翻译和可核验约束组合评分：
\begin{equation}
 S(x,y)=\lambda_a\log p_{\mathrm{asr}}(x\mid\bm{X})
        +\lambda_m\log p_{\mathrm{mt}}(y\mid x)+r(x,y).
 \label{eq:nlp-speech-nbest}
\end{equation}
$r$可检查日期、实体和格式的一致性，权重需在开发集按相同候选预算校准；不同模型的对数概率受分词与长度影响，不应直接把数值大小当成可比置信度。译文生成参数来自平行数据，不因为多传几个识别候选就自动重新训练。

若有完整的概率模型，可以定义$p(y\mid\bm{X})=\sum_xp_{\mathrm{mt}}(y\mid x)p_{\mathrm{asr}}(x\mid\bm{X})$。实际$n$最佳列表只覆盖部分$x$，而逐候选取最佳译文再重排通常是在近似最大化路径对分数，不是对相同译文完整边缘化。两个源候选译成相同文字时，也要先说明合并规则。

教学候选“十八日”和“十七日”的识别概率为$.55,.45$，各自最佳译文条件概率为$.60,.90$；直接乘积是$.33,.405$，后者获胜，说明语言生成偏好可能压过较好的声学候选。加入规范值检查可以拒绝源译不一致，却不能在未知源语真值时替我们决定哪个日期确实说过。比较应同时报告候选里存在有效译文的比例、最终选择正确的比例与额外解码成本。

\subsubsection{增量级联与等待策略}
\label{sec:nlp-speech-incremental-cascade}

用户没有说完时，级联各环节不能假定完整源语已到达。声音块到达识别器后，只有已提交的稳定源语前缀才能成为不可撤回翻译的依据；暂定文字可用于预计算，但必须带版本并允许废弃。固定声音块、文本wait-$k$和自适应等待控制的是不同位置：文本wait-$k$在读入$k$个源文字单位后开始写，再随目标步推进读入，不能把一个源词等同于固定毫秒的声音。

例如用户说“校外来宾可以直接参加……吗？不，须先报名”。过早把“可以直接参加”翻译并播放，会把尚未完成的询问误作许可。翻译训练可用带规定读写掩码的平行数据，或在既有模型外加入提交策略；推理记录每次可见声音终点、识别确认前缀、目标暂定尾部及已提交输出。模型预测未来内容不等于已经获得源声音证据。

目标文字尚未展示时可以重算；已展示文字能否修订由产品协议决定；已播放的声音只能明确更正，不能删除。故延迟至少拆成声音收集与识别等待、翻译等待、合成生成及播放队列等待。假设用户于墙钟3.0\,s结束，最终源文3.4\,s提交，译文3.7\,s确定，首段声音4.0\,s开始播放，则结束到首段播放为1.0\,s，而非翻译阶段的0.3\,s。这是计时演算，不是系统实测。

\subsection{直接转换}
\label{sec:nlp-speech-direct-translation}

直接路线由源语声学编码接目标语言输出，不必先调用外部识别器取得源语转写。训练可使用声音—译文对，也可通过ASR辅助、文本翻译迁移和联合训练利用其他监督。省去外部转写接口并不排除模型内部有文字阶段，也不说明声学错误已经消失。与级联比较时，应固定语言方向、数据资源、可见上下文和计算预算。

\subsubsection{文本与声音单位生成：SeamlessM4T}
\label{sec:nlp-speech-seamless}

SeamlessM4T把声音或文字输入接到目标文字解码器，需要声音时再生成离散声音单位，并由单位声码器恢复波形。它采用UnitY式两阶段生成：第一阶段输出目标语言文字，第二阶段读取该文字的解码表示，预测目标声音单位。这里内部文字是目标译文，不是强制要求先交付源语转写；所以“直接语音翻译”与“内部完全没有文字”不是同义词。原论文与完整作者信息列于文献说明。

声音编码器先从未标注录音预训练，文字翻译部件从平行文本学习，再用声音—译文、识别与其他相应样本联合训练目标文字。声音单位则由另一声音表示模型提取目标录音特征，再聚类成离散索引；索引词表来自声音数据，不等于文字词表，也不是人工逐音素标注。文字到单位模型以目标语言转写与对应声音单位训练，单位声码器另用声音单位和真实波形配对学习。伪标注和挖掘配对可以补充监督，仍有错配与领域变化风险。

使用时，声音编码器接收源语通知及目标语言条件，文字解码器束搜索形成译文；将选中的译文及其解码表示送入文字到单位模型，再搜索单位序列，补足单位时长并经条件声码器形成目标波形。原模型的两阶段最优候选选择不等于对所有文字与单位组合做全局最优搜索。若要求仅输出文字，可以在第一阶段结束；要求声音就必须继续检查第二阶段和波形。

对“海棠楼201”，译文可在既定译名表下使用“Haitang Building, Room 201”，但译名表不能凭空新增楼名。若第一阶段保留201、第二阶段声音读成210，文字评价会漏掉错误；若文字和声音都正确却新增“无需报名”，仍违反通知原意。每个语言方向、声音输入条件及文字或声音输出都须独立评价，不能从一组方向的平均分推出所有组合的保证。

\subsubsection{单调推进与流式译出：SeamlessStreaming}
\label{sec:nlp-speech-seamless-streaming}

完整源语未到达时，直接翻译模型也要决定继续读入还是现在输出。SeamlessStreaming在SeamlessM4T v2上加入高效单调多头注意力（efficient monotonic multihead attention, EMMA），用每个注意力头的读写决策约束可用源语位置；随后由UnitY2式非自回归文字到单位部件生成对应声音片段。它与上一小节原始SeamlessM4T的单位生成方式有所不同，不能因名称相近而混用实现。

对一个头，令$p_{i,j}$为已生成目标前缀时，在源语位置$j$允许生成第$i$个目标单位的概率，由上一解码状态和当前源语表示经可学习映射及sigmoid得到。单调对齐概率可写成
\begin{equation}
 \alpha_{i,j}=p_{i,j}
      \sum_{k=1}^{j}\alpha_{i-1,k}
         \prod_{\ell=k}^{j-1}(1-p_{i,\ell}).
 \label{eq:nlp-speech-emma}
\end{equation}
从前一个边界$k$走到$j$，意味着沿途暂不输出，到$j$才选择输出。训练对这些可能边界求期望，使用允许回看已读前缀的注意力；EMMA给出稳定的并行估计，避免直接使用易产生数值问题的比例式计算。翻译负对数概率还加上期望延迟与对齐方差正则，否则模型可能学成等全部声音到齐才输出的退化离线策略。原文分阶段训练同步文字解码器和文字到单位部件，具体来源列于文献说明。

推理每次新声音块到达，原文实现重新计算当前已见声音的编码表示，检查所有头的决策。只有最保守的头也达到阈值时才写目标词元，否则继续读入。教学设定三个头的概率为$(.7,.6,.3)$、阈值$.5$，此时等待；更多声音到达后变为$(.8,.7,.6)$，才允许输出。这些概率来自构造，说明多头决策，不表示已测出的响应时刻。

新生成的文字及解码状态进一步变成对齐的声音单位；单位块达到最小长度才交给声码器，所以“允许写一个词”不等于“立即播放一个词”。对句末才确认的日期或否定，应同时检查是否过早形成错误译文、是否积累了过长声音块，以及已播部分怎样纠正。源语位置单调推进不要求逐词直译，重排序需要的未来信息仍可能迫使系统等待。

延迟有两种不同口径。以“生成第$i$个目标单位时已经读入多少秒源声音”衡量的滞后，描述信息等待；以墙钟测得的首段播放、最后声音播放结束，还包含编码重算、解码、网络与队列成本。原文语音输出的结束偏移比较源声音结束和目标声音结束，本书实际服务评价还应记录对应墙钟事件。声音表达、语速和停顿的保持属于额外目标，不能由低滞后或模型名称直接推出。

\section{实时口语对话}
\label{sec:nlp-speech-realtime-dialogue}

语音输入输出单独成立后，完整对话仍要协调何时听、何时说、当前要求是什么、工具是否真的完成动作。\term{全双工}指双方可以同时接收和发送声音的交互条件，不等于每次交叠都应继续播放。一个支持持续听说的模型可能不擅长业务状态更新，一个串联识别、文本对话与合成的系统也可能实现合理打断；必须检查实际声音事件与状态记录。

\subsection{输入流、轮次与打断}
\label{sec:nlp-speech-turns}

“帮我订……明天下午”的停顿可能是思考，而非话轮结束。系统播报时用户说“等等”，可能要求接管或更正；短促“嗯”则可能仅表示在听。任务输入应保留双方声音时间轴、播放状态和可用历史，监督分别标注声音活动、轮次边界与交流事件，不能把它们压成一个“说完了”标签。

\subsubsection{VAD与轮次结束检测}
\label{sec:nlp-speech-vad}

\term{语音活动检测}（voice activity detection, VAD）判断某段是否存在语音。一个教学能量基线对短窗计算$E_i=w^{-1}\sum_nx_{ih+n}^2$，超过启动阈值$\theta_{\mathrm{on}}$时进入活动状态，连续若干窗低于结束阈值$\theta_{\mathrm{off}}$才退出，并令$\theta_{\mathrm{on}}>\theta_{\mathrm{off}}$以减少来回抖动。阈值从噪声校准和开发录音选择；神经VAD则以有声、无声帧标注训练分类器，再加同类平滑。能量高可能是噪声，能量低也可能是轻声，因此两者都需要声音条件测试。

轮次结束检测还要知道语句和交流作用是否完整。可设“听取—候选结束—等待—确认结束”状态机：VAD转静音时进入候选结束，结合静音长度、韵律下降、语义完整度和历史判断；若用户恢复发声，则回到听取并撤销尚未提交的结束判断。教学阈值设为400\,ms，用户在“帮我订”后停顿300\,ms再说“明天下午”，不应仅因出现静音就开始正式回答；即使停顿达到500\,ms，明显缺少预约对象也可以触发澄清而非执行。

对VAD评估活动区间，对结束检测评估过早截断和结束等待；用户真正交出话轮的标注需要独立协议。在线阈值不能用完整未来录音偷偷修正后再计为低延迟。系统播放自己的声音还可能回到麦克风，回声处理或双通道信息决定VAD实际接收什么，也应写入测试条件。

\subsubsection{未来声音活动预测：VAP}
\label{sec:nlp-speech-vap}

静音长度只是已经发生的证据，轮次控制还可预测接下来谁会说话。未来声音活动预测（voice activity projection, VAP）用已见声音与双方活动历史预测未来两人的活动模式。原论文把未来2秒划为200、400、600、800\,ms四个连续区间，每人四位，以区间内活动比例是否超过阈值决定0或1；两人合计8位，对应$2^8=256$种联合状态。完整来源见文献说明。

训练标签由录音后续的活动标记构造，模型输入只见当前及过去，不能泄漏未来窗口。离散VAP用交叉熵学习256类分布，因而能表示两个说话者在各区间之间的依赖。推理将与某事件相符的状态概率求和，必要时在比较的事件子集内归一化，得到保持或移交话轮、短回应等信号；比逐区间独立阈值化更明确地保留了联合状态意义。

例如教学简化状态“A在前两区间活动、之后不活动；B在后两区间活动”是移交的一种候选，但真实事件定义还要规定之前是谁在说、是否互相静音及之后持续多久。若符合移交、保持和其他模式的总质量为$.55,.35,.10$，二选一条件下移交概率为$.55/(.55+.35)=.611$，而无条件移交质量仍为$.55$。两种数值不能混用。

短回应和长插话都可能以同一声“嗯”开始，VAP应在固定可见窗口上预测并与后续实际活动比较，而不是事后按长度重写预测记录。该信号可以辅助是否让出话轮，但不能证明预约信息齐全，也不能证明用户授权了工具操作。

\subsection{增量理解与响应选择}
\label{sec:nlp-speech-incremental-understanding}

实时识别会把“明天”修成“后天”，连续假设不是不断追加的文字日志。每次局部转写或声音语义应带版本、来源区间以及已确认前缀和暂定尾部。状态更新按新增、替换和撤回处理：版本1暂记4月18日，版本2听到“不，后天”后改为4月19日，旧日期及其派生的未执行动作都要失效。历史原话可以保留供追溯，当前有效值只能有符合本轮更正的一份。

\subsubsection{状态驱动的响应门控}
\label{sec:nlp-speech-response-gate}

响应门控在状态和事件条件下选择等待、简短确认、澄清或正式回答。一个规则实现检查语义是否完整、关键字段是否确认、用户是否交出话轮、工具是否返回以及剩余预算。用户仍在更正时继续听；缺日期时询问日期；参数齐全但授权缺失时请求确认；工具未返回时至多说“正在查询”，只有返回提供结果依据后才回答结果。规则来源是业务前置条件和交互协议，也可用按同一动作集标注的历史训练行为分类器。

这些分支应改变明确状态。对“明天……不，后天”，第一次仅创建暂定日期，第二次替换日期并取消依赖旧版本的待发请求；对已发送查询，则等待其返回时检查版本，不把旧结果写回新状态。行为生成模型可以提出动作，但必须经过同样的字段、权限和结果条件检查。短应答回馈（backchannel）如“嗯，我在听”只表达接收，不自动意味着业务许可或完成。

教学门控若看到“日期已确认、人数缺失、轮次结束”，应输出澄清人数并进入等待用户状态；若看到“字段完整、查询在途”，应等待而非伪造可预约时段。评价不仅看动作标签，还看是否错误抢话、把暂定值提交给工具，以及澄清是否真正补回缺失信息。规则阈值和学习模型均须在留出的完整对话中检查。

\subsubsection{平行声音流与持续听说：Moshi}
\label{sec:nlp-speech-moshi}

Moshi用用户和系统两个平行声音流表示持续交互，每个流在共同时间网格上由神经音频编解码器转换成多层声音词元，系统文字流与其将说出的声音对齐。模型沿时间建模历史，再在当前时间步内按文字、语义性声音信息与其余声学词元的顺序生成。原文将系统文字作为声音生成的先行条件称为Inner Monologue；它不是给用户发言先做一份外部实时转写。完整来源列于文献说明。

训练以录音、对应文字对齐和对话声音流形成词元监督，对应位置的交叉熵更新模型；音频编解码器另学习从波形到离散表示及恢复波形。原实现声音帧率为12.5\,Hz，即每步80\,ms，词元流间还有明确的延迟安排。帧率只规定表示节拍，不等于80\,ms的端到端回复延迟，计算、输入积累和播放仍会增加等待。

推理时用户词元来自实际收到的用户声音，不采用模型猜测出的用户词元；系统文字和系统声音词元由模型生成，编解码器恢复系统声音。即使系统正在说话，用户流仍能继续输入，所以表示中可以保留交叠，也可以输出近静音。与识别—文本对话—合成流水线相比，差别在于持续声音如何进入上下文、声音与文字如何共同预测，而不是自动取消所有外部状态管理。

对系统刚说“按明天……”时用户更正“不是，后天”，要检查用户声音是否进入模型、是否停止旧回复、后续声音是否改用正确日期，以及业务状态是否同步。联合声音生成不自动让工具结果可见，也不保证数据库日期已更新。若应用接入工具，仍须提供可检查的任务状态接口，不能以模型自然地接话代替执行证据。

\subsection{声音回复与行为协调}
\label{sec:nlp-speech-coordination}

实时回复应把文字或声音片段分成待生成、已生成未播、已播三类。每个片段保存任务版本、依赖字段与播放位置；新的更正或工具结果使旧版本失效时，取消过时生成任务并清空未播队列。若用户确实试图接管，停止或衰减当前播放，记录已经播出的范围，再处理更正。短回应与接管需区别，否则用户每说一次“嗯”，系统都会丢弃有用回复。

工具请求应附请求标识、任务版本及完成状态。取消本地等待不一定取消服务端写操作；收到迟到结果时应核查实际执行记录，不可直接重发写请求。幂等、查询执行结果和失败恢复由第\ref{chap:nlp-solving-and-execution}章的执行闭环处理，声音层需要保留这些证据并决定什么时候播报。工具未返回可以给有依据的进度确认，不能说“已经改好了”。

构造一条墙钟记录：用户于3.0\,s说完，识别3.4\,s提交，语义与查询请求在3.5\,s确认，系统3.8\,s开始播放“正在查询”，工具4.6\,s返回，结果声音5.0\,s开始播放。结束到首段声音为0.8\,s，结束到结果声音为2.0\,s；若业务完成依据就是4.6\,s的已核验返回，从输入开始计的完成时间为4.6\,s。简短确认改善了反馈等待，但没有把业务完成提前到3.8\,s。

再设用户于4.0\,s插话要求改日期，系统4.15\,s停止播放，打断停止延迟为150\,ms。此时4.6\,s返回的旧日期查询只能保存为旧版本结果，不能播成对新日期的回答。已经播出的错误日期须明确更正，不能从内部历史中删除后假装用户没听到。文字接增量TTS与联合声音词元生成都需要这套播放和状态协调。

$\tau$-Voice把业务结果与语音交互质量放在同一评价框架内：按环境终态检查任务成功，并记录回应、让出话轮、打断和选择性处理短回应等行为。它采用离散模拟声音时间推进，与墙钟计算解耦；用户模拟器还直接接收系统转写，而非再次识别系统声音。其公开协议因此可用于组织故障测试，但不独自验证实际声音生成质量，也不能把模拟时间直接当成真实部署的墙钟延迟\scite{nlp-intro-ray2026}

实际测试应保存用户结束、稳定识别、首段播放、结果播放、打断停止和业务完成事件，并记录首段等待、总时间、错误抢话、重复操作和失败恢复。口音、噪声、多语、交叠与接口失败应进入完整交互测试；局部ASR和TTS分数帮助定位问题，却不能取代完整任务成功与用户实际听到的结果。只有这些事件共享明确时间基准，表\ref{tab:nlp-speech-contracts}中的时间比较才成立。

\subsection*{文献说明}
\label{sec:nlp-speech-literature}

下列较长或落点相邻的条目集中列出，以便对应正文所述的理解、合成、翻译与持续听说实现。SLURP与VAP分别支持口语语义评价和未来活动预测；Tacotron 2的来源支持字符到频谱和条件WaveNet；两篇Seamless论文分别支持原始两阶段多任务生成与后续流式策略；Moshi来源支持平行声音流及声文时间组织。本文的中文通知、数值格点、状态机和业务核验例子不属于这些论文的实验。

\begin{fullwidth}
  \raggedright
  \fullcite{nlp-speech-bastianelli2020}\par
  \fullcite{nlp-speech-ekstedt2022}\par
  \fullcite{nlp-intro-shen2018}\par
  \fullcite{nlp-intro-seamless2023}\par
  \fullcite{nlp-speech-seamless2023streaming}\par
  \fullcite{defossez2024moshi}
\end{fullwidth}

\section{本章小结}
\label{sec:nlp-speech-summary}

从声音取得信息，需要同时说明听见了什么、符号怎样对应时间，以及哪些结果已经稳定。CTC对可折叠路径求和，RNN-T把输出历史加入时间—输出格点，注意力生成则按前缀读取声音；不同的概率与对齐机制都不能单独证明流式能力。声音理解进一步恢复有效日期、数量、更正与交流作用，不能以少量文字错误掩盖关键字段失败。

声音输出把保持内容与选择表达分开。读法规范、发音、时长和韵律、频谱与波形可以由多个模块或联合模型承担，但需要逐项定位错误，并在最终声音中核验。翻译又增加跨语言意义保持，目标文字正确并不保证目标声音正确；直接路线省去某些外部接口，也不省去监督来源和结果检查。

章首的日期询问要求系统依据声音和历史理解有效日期、适时给出可听回应，并在涉及业务动作时等待真实完成证据。记录局部文字、模型时间网格与墙钟事件，才能分别核验准确性、等待和修订代价。图像、页面和音画联合证据还需要视觉输入与跨模态核验，由后续图像部分展开。
